\scanpage{47}
$M$ motif sur $k$ \inserted{alg. clos}, $k$ isom à un sous-corps de $\mathbb C$

Pour toute place $v$ de $k$, on a $M_v^{\mathbb Z}$ qui est
un sous-réseau de $\prod_{l\ \mathrm{fini}}M(l)$. On veut établir
les relations entre les $M_v^{\mathbb Z}$ pour diverses
valeurs $v_1\ldots v_n$ de $v$. Choisissons un
bon foncteur de référence (p. ex. défini par $v_0$)
$\unsure{\xi|\zeta}\ M\unsure{\sim|\mapsto|\rightsquigarrow} M_{\xi_0}\quad(\xi_0\in T(\mathbb Q))$ avec \deleted{des} isom \unsure{fonct.|fonctoriels|foncto}
$\xi_{0\deleted{\unsure{\mathbb Q_l}}}(\deleted{\unsure{H}}M)_{\mathbb Q_l}\simeq M_l\otimes_{\mathbb Z_l}\mathbb Q_l$. Alors $v_1,\ldots,v_n$ définissent des
fibrés principaux homogènes
\[
P_i=\underline{\mathrm{Isom}}(\xi,\xi_{v_i})
\]
sous $G_{M\xi_0}^{\unsure{k|f|\xi}}$ (sur $\mathbb Q$) et on a un
point de $(\prod P_i)(\unsure{\widehat{\overline{\mathbb Q}}|\widehat{\mathbb Q}})$, où $\widehat{\mathbb Q}=\prod_{l\ \mathrm{fini}}\mathbb Q_l$,
d'où un pt de $[(\prod P_i)/G](\widehat{\mathbb Q})$.
% slanted
\margin{$G$ opère diagonalement}

On peut aussi considérer que l'on a un
point canonique dans \deleted{$\mathrm{Gl}(n)\times\ldots\times\mathrm{Gl}(n)$}
\deleted{\unsure{$\mathrm{Gl}$}}
\[
\begin{aligned}
&\deleted{\backslash[\mathrm{Gl}(n)\times\ldots\times\mathrm{Gl}(n)](\widehat{\mathbb Q})/\mathrm{Gl}(n)(\widehat{\mathbb Q})}\\
&\deleted{(\mathrm{Gl}(n)\times\ldots\times\mathrm{Gl}(n))/\mathrm{Gl}(n)}\\
&\mathrm{Gl}(m,\mathbb Z)\times\ldots\mathrm{Gl}(m,\mathbb Z)\backslash
\mathrm{Gl}(m,\widehat{\mathbb Q})\times\ldots\times\mathrm{Gl}(m,\widehat{\mathbb Q})/
\mathrm{Gl}(m,\widehat{\mathbb Q})\qquad m=\text{rang }M
\end{aligned}
\]

Les considérations qui précèdent montrent que
le degré de transcendance du \unsure{pt|pts} en question
% archive sheet number: 46
